\chapter{Erythema Nodosum}

\section*{Definition and Overview}
Erythema nodosum is a skin condition in which painful, red, tender lumps develop, most commonly on the shins. The lumps are a form of inflammation of the fat layer beneath the skin, called panniculitis. Erythema nodosum is not a disease in itself but a reaction that can be triggered by many conditions, including infections, medicines, pregnancy, and inflammatory diseases such as sarcoidosis. It usually resolves on its own within a few weeks, but it is important to find and treat the underlying cause.

\section*{Epidemiology}
Erythema nodosum is the most common form of panniculitis and affects people of all ages, most commonly young and middle-aged adults, and women more often than men. It can occur at any time but is more common in spring and autumn in some regions. In many cases no cause is found, while in others a specific trigger is identified after investigation.

\section*{Aetiology and Pathophysiology}
Erythema nodosum is a hypersensitivity reaction in the fat layer of the skin.
\begin{itemize}
    \item \textbf{Infections:} streptococcal throat infections are the most common trigger in children; other causes include tuberculosis, fungal infections, and viral illnesses
    \item \textbf{Sarcoidosis:} an inflammatory disease that is a common underlying cause, particularly in adults
    \item \textbf{Inflammatory bowel disease:} Crohn's disease and ulcerative colitis
    \item \textbf{Medicines:} some antibiotics, oral contraceptives, and other drugs
    \item \textbf{Pregnancy:} erythema nodosum can occur during pregnancy
    \item \textbf{Mechanism:} an immune reaction in the fat layer causes inflammation, producing the characteristic tender nodules
    \item \textbf{The nodules:} they begin red and swollen, then fade to purple and yellow-brown as they heal, without scarring or ulceration
\end{itemize}

\section*{Clinical Features}
The skin findings are characteristic.
\begin{itemize}
    \item Painful, red, tender lumps on the shins, often on both legs
    \item The lumps may also appear on the thighs, arms, or elsewhere
    \item The nodules are warm and slightly raised, usually 1 to 5 centimetres across
    \item They fade to bruise-like colours over one to three weeks
    \item Fever, fatigue, and joint pain may accompany the rash
    \item Swelling of the ankles in some cases
    \item The condition does not scar or ulcerate
\end{itemize}

\section*{Differential Diagnosis}
Other skin conditions can look similar.
\begin{itemize}
    \item Cellulitis, a bacterial skin infection with spreading redness and fever
    \item Insect bites
    \item Vasculitis, inflammation of the blood vessels
    \item Thrombophlebitis, inflammation of a vein
    \item Other forms of panniculitis
    \item Bruising and trauma
\end{itemize}

\section*{When to Seek Medical Care}
Anyone with painful lumps on the legs should be assessed by a GP, both to confirm the diagnosis and to look for an underlying cause. Medical care is needed for the underlying conditions that can trigger erythema nodosum.

\textbf{Call 000 (or local emergency services) for:}
\begin{itemize}
    \item Fever with a rapidly spreading, painful, red area of skin, which may indicate cellulitis
    \item Difficulty breathing, which may indicate sarcoidosis involving the lungs
    \item Chest pain or breathlessness
    \item Collapse or severe illness
\end{itemize}

\section*{Diagnosis and Investigations}
Investigations aim to find the underlying cause.
\begin{itemize}
    \item \textbf{History and examination:} the appearance of the rash and any recent infection, medicines, or pregnancy
    \item \textbf{Throat swab and blood tests:} for streptococcal infection and inflammatory markers
    \item \textbf{Chest X-ray:} to look for sarcoidosis and tuberculosis
    \item \textbf{Other tests:} for inflammatory bowel disease, fungal infection, and other causes, guided by the history
    \item \textbf{Skin biopsy:} rarely needed, when the diagnosis is uncertain
\end{itemize}

\section*{Management}
Treatment of erythema nodosum is usually about comfort and treating the cause.
\begin{itemize}
    \item \textbf{Pain relief:} anti-inflammatory medicines such as ibuprofen
    \item \textbf{Rest and elevation:} resting the legs and elevating them reduces discomfort
    \item \textbf{Compression:} supportive stockings can help
    \item \textbf{Treating the cause:} antibiotics for streptococcal infection, and treatment of any underlying disease
    \item \textbf{Stopping the trigger:} ceasing a causative medicine under medical advice
    \item \textbf{Severe cases:} corticosteroids or other treatment for persistent or severe disease
    \item \textbf{Follow-up:} the rash is monitored as it heals
\end{itemize}

\section*{Complications}
\begin{itemize}
    \item The pain and discomfort of the nodules
    \item Recurrence, which occurs in some people
    \item The complications of the underlying disease, such as sarcoidosis or inflammatory bowel disease
    \item Joint pain and swelling, which can accompany the condition
\end{itemize}

\section*{Prognosis and Outcomes}
Erythema nodosum usually resolves on its own within one to six weeks, even without treatment, and it does not leave scars. The outlook depends on the underlying cause: when a specific trigger is found and treated, the condition settles completely, while in sarcoidosis and inflammatory bowel disease it may recur with the underlying illness.

\section*{Prevention and Self-Care}
\begin{itemize}
    \item See a doctor for any new painful lumps on the legs
    \item Complete the full course of treatment for any underlying infection
    \item Rest and elevate the legs while the rash is active
    \item Use anti-inflammatory pain relief as advised
    \item Attend follow-up to complete the search for a cause
    \item Report any recurrence or new symptoms
\end{itemize}

\section*{Sources and Further Reading}
The following reputable sources provide further information for healthcare students and patients seeking reliable, current advice about erythema nodosum.
\begin{itemize}
    \item DermNet NZ and DermNet Australia. \textit{Erythema nodosum.} https://dermnetnz.org/
    \item Healthdirect Australia. \textit{Erythema nodosum.} https://www.healthdirect.gov.au/
    \item Sarcoidosis Australia. \textit{Sarcoidosis and skin symptoms.} https://www.sarcoidosis.org.au/
    \item NHS. \textit{Erythema nodosum.} https://www.nhs.uk/conditions/erythema-nodosum/
    \item Mayo Clinic. \textit{Erythema nodosum.} https://www.mayoclinic.org/diseases-conditions/erythema-nodosum/
    \item MSD Manual. \textit{Erythema nodosum.} https://www.msdmanuals.com/
\end{itemize}
